\subsection{Making Review Binding}
\label{sec:making-binding}

Review becomes binding when an adverse finding changes the institution's
ability to continue without requiring the institution to agree with the
finding. Expertise, independence, publication, and reputational pressure
may improve review. None of them supplies that consequence.

Human-subjects research provides a working example. An institutional
review board does not ordinarily fine or prosecute a researcher.
Instead, an institution promises, through a federalwide assurance, that
research within the assurance's scope will receive and comply with
review. A federal office can suspend the assurance when the institution
breaks that promise, placing the institution's eligibility for federal
research funding at risk rather than merely censuring the individual
laboratory \autocite{ohrp2018fwa}. The ruling binds because another
institution controls something the researcher needs in order to
continue.

The transferable part is the chain of authority, not federal funding in
particular. A jurisdiction can make review-board clearance a condition
of the license under which research on candidate sentient systems is
conducted. Suspension of clearance would then suspend legal authority to
run the protocol, and continued operation would expose the institution
to the ordinary enforcement available for unlicensed research. The
board would not need to punish the institution itself. Its finding would
activate a consequence administered by a separate licensing authority.

Each link matters. The board must have jurisdiction before the disputed
protocol begins, access to evidence the operator cannot withhold, and
authority to suspend clearance. The licensing body must be obliged to
act on that suspension rather than invited to reconsider it. An appeal
may run to a separate body, but the operator cannot decide for itself
that the appeal permits the work to continue. If another discretionary
decision by the reviewed institution is required anywhere in the chain,
the ruling is advisory at that point.

\standing{The rest of this section applies the chain to research on a
candidate moral patient, and assumes the route.}
\runin{What the board applies}
The substantive standard is animal-research ethics' Three Rs \autocite{russell1959principles}. A binding board is what makes those three questions something a protocol has to answer before it runs. In practice, applying the Three Rs to a candidate sentient system means asking, in order: can a simulation or a non-sentient model answer the same question; if not, what is the minimum exposure that still produces a valid result; and given that minimum, what escalation triggers stop the study before harm compounds. A board with teeth can require a negative answer to the first question be documented, not merely asserted, before it will consider the second.

The Three Rs' own logic — defaulting toward the least invasive option available — argues for a specific evidentiary rule at every stage of this apparatus: a system's default experimental status stays the more protective one until the case for the less protective status is made, not the reverse. That default has a cost, and neither error it produces is free: it will sometimes extend protection to a system that turns out not to have needed it, and sometimes, before a threshold is met, withhold protection from one that did. The Three Rs answer only what a board does once a subject is inside its scope.
